\subsection{Basic definitions}

\begin{definition}
\label{def.perm.perm}Let $X$ be a set. \medskip

\textbf{(a)} A \emph{permutation} of $X$ means a bijection from $X$ to $X$.
\medskip

\textbf{(b)} It is known that the set of all permutations of $X$ is a group
under composition. This group is called the \emph{symmetric group} of $X$, and
is denoted by $S_{X}$. Its neutral element is the identity map
$\operatorname*{id}\nolimits_{X}:X\rightarrow X$. Its size is $\left\vert
X\right\vert !$ when $X$ is finite.

(Alternative notations for $S_{X}$ include $\operatorname*{Sym}\left(
X\right)  $ and $\Sigma_{X}$ and $\mathfrak{S}_{X}$ and $\mathcal{S}_{X}$.)
\medskip

\textbf{(c)} As usual in group theory, we will write $\alpha\beta$ for the
composition $\alpha\circ\beta$ when $\alpha,\beta\in S_{X}$. This is the map
that sends each $x\in X$ to $\alpha\left(  \beta\left(  x\right)  \right)  $.
\medskip

\textbf{(d)} If $\alpha\in S_{X}$ and $i\in\mathbb{Z}$, then $\alpha^{i}$
shall denote the $i$-th power of $\alpha$ in the group $S_{X}$. Note that
$\alpha^{i}=\underbrace{\alpha\circ\alpha\circ\cdots\circ\alpha}_{i\text{
times}}$ if $i\geq0$. Also, $\alpha^{0}=\operatorname*{id}\nolimits_{X}$.
Also, $\alpha^{-1}$ is the inverse of $\alpha$ in the group $S_{X}$, that is,
the inverse of the map $\alpha$.
\end{definition}

\begin{definition}
\label{def.perm.Sn-iven}Let $n\in\mathbb{Z}$. Then, $\left[  n\right]  $ shall
mean the set $\left\{  1,2,\ldots,n\right\}  $. This is an $n$-element set if
$n\geq0$, and is an empty set if $n\leq0$.

The symmetric group $S_{\left[  n\right]  }$ (consisting of all permutations
of $\left[  n\right]  $) will be denoted $S_{n}$ and called the $n$\emph{-th
symmetric group}. Its size is $n!$ (when $n\geq0$).
\end{definition}

For instance, $S_{3}$ is the group of all $6$ permutations of the set $\left[
3\right]  =\left\{  1,2,3\right\}  $.

If two sets $X$ and $Y$ are in bijection, then their symmetric groups $S_{X}$
and $S_{Y}$ are isomorphic. Intuitively, this is clear (just think of $Y$ as a
\textquotedblleft copy\textquotedblright\ of $X$ with all elements relabelled,
and use this to reinterpret each permutation of $X$ as a permutation of $Y$).
We can formalize this as the following proposition:

\begin{proposition}
\label{prop.perm.Sf}Let $X$ and $Y$ be two sets, and let $f:X\rightarrow Y$ be
a bijection. Then, for each permutation $\sigma$ of $X$, the map $f\circ
\sigma\circ f^{-1}:Y\rightarrow Y$ is a permutation of $Y$. Furthermore, the
map%
\begin{align*}
S_{f}:S_{X}  &  \rightarrow S_{Y},\\
\sigma &  \mapsto f\circ\sigma\circ f^{-1}%
\end{align*}
is a group isomorphism; thus, we obtain $S_{X}\cong S_{Y}$.
\end{proposition}

\begin{proof}
Easy and LTTR.
\end{proof}

Because of Proposition \ref{prop.perm.Sf}, if you want to understand the
symmetric groups of finite sets, you only need to understand $S_{n}$ for all
$n\in\mathbb{N}$ (because if $X$ is a finite set of size $n$, then there is a
bijection $f:X\rightarrow\left[  n\right]  $ and therefore a group isomorphism
$S_{f}:S_{X}\rightarrow S_{\left[  n\right]  }$). Thus, we will focus mostly
on $S_{n}$ in this chapter.

\begin{remark}
If $Y=X$ in Proposition \ref{prop.perm.Sf}, then the group isomorphism $S_{f}$
is conjugation by $f$ in the group $S_{X}$.
\end{remark}

Next, let us define three ways to represent a permutation:

\begin{definition}
\label{def.perm.notations}Let $n\in\mathbb{N}$ and $\sigma\in S_{n}$. We
introduce three notations for $\sigma$: \medskip

\textbf{(a)} A \emph{two-line notation} of $\sigma$ means a $2\times n$-array
\[
\left(
\begin{array}
[c]{cccc}%
p_{1} & p_{2} & \cdots & p_{n}\\
\sigma\left(  p_{1}\right)  & \sigma\left(  p_{2}\right)  & \cdots &
\sigma\left(  p_{n}\right)
\end{array}
\right)  ,
\]
where the entries $p_{1},p_{2},\ldots,p_{n}$ of the top row are the $n$
elements of $\left[  n\right]  $ in some order. Note that this is a standard
notation for any kind of map from a finite set. Commonly, we pick $p_{i}=i$,
so we get the array
\[
\left(
\begin{array}
[c]{cccc}%
1 & 2 & \cdots & n\\
\sigma\left(  1\right)  & \sigma\left(  2\right)  & \cdots & \sigma\left(
n\right)
\end{array}
\right)  .
\]


\textbf{(b)} The \emph{one-line notation} (short, \emph{OLN}) of $\sigma$
means the $n$-tuple $\left(  \sigma\left(  1\right)  ,\sigma\left(  2\right)
,\ldots,\sigma\left(  n\right)  \right)  $.

It is common to omit the commas and the parentheses when writing down the OLN
of $\sigma$. Thus, one simply writes $\sigma\left(  1\right)  \ \sigma\left(
2\right)  \ \cdots\ \sigma\left(  n\right)  $ instead of $\left(
\sigma\left(  1\right)  ,\sigma\left(  2\right)  ,\ldots,\sigma\left(
n\right)  \right)  $. Note that this omission can make the notation ambiguous
if some of the $\sigma\left(  i\right)  $ have more than one digit (for
example, the OLN $1112345678910$ can mean two different permutations of
$\left[  11\right]  $, depending on whether you read the \textquotedblleft%
$111$\textquotedblright\ part as \textquotedblleft$1,11$\textquotedblright\ or
as \textquotedblleft$11,1$\textquotedblright). However, if $n\leq10$, then
this ambiguity does not occur, and the notation is unproblematic (even without
commas and parentheses). \medskip

\textbf{(c)} The \emph{cycle digraph} of $\sigma$ is defined (informally) as follows:

\begin{itemize}
\item For each $i\in\left[  n\right]  $, draw a point (\textquotedblleft
node\textquotedblright) labelled $i$.

\item For each $i\in\left[  n\right]  $, draw an arrow (\textquotedblleft
arc\textquotedblright) from the node labelled $i$ to the node labelled
$\sigma\left(  i\right)  $.
\end{itemize}

The resulting picture is called the cycle digraph of $\sigma$.

Using the concept of \emph{digraphs} (= directed graphs), this definition can
be restated formally as follows: The \emph{cycle digraph} of $\sigma$ is the
directed graph with vertices $1,2,\ldots,n$ and arcs $i\rightarrow
\sigma\left(  i\right)  $ for all $i\in\left[  n\right]  $.
\end{definition}

\begin{example}
Let $\sigma:\left[  4\right]  \rightarrow\left[  4\right]  $ be the map that
sends the elements $1,2,3,4$ to $2,4,3,1$, respectively. Then, $\sigma$ is a
bijection, thus a permutation of $\left[  4\right]  $. \medskip

\textbf{(a)} A two-line notation of $\sigma$ is $\left(
\begin{array}
[c]{cccc}%
1 & 2 & 3 & 4\\
2 & 4 & 3 & 1
\end{array}
\right)  $. Another is $\left(
\begin{array}
[c]{cccc}%
3 & 1 & 4 & 2\\
3 & 2 & 1 & 4
\end{array}
\right)  $. Another is $\left(
\begin{array}
[c]{cccc}%
4 & 1 & 3 & 2\\
1 & 2 & 3 & 4
\end{array}
\right)  $. There are $24$ two-line notations of $\sigma$ in total, since we
can freely choose the order in which the elements of $\left[  4\right]  $
appear in the top row. \medskip

\textbf{(b)} The one-line notation of $\sigma$ is $\left(  2,4,3,1\right)  $.
Omitting the commas and the parentheses, we can rewrite this as $2431$.
\medskip

\textbf{(c)} One way to draw the cycle digraph of $\sigma$ is
\[%
% [tikz figure removed]
%BeginExpansion
% [tikz figure removed]%
%EndExpansion
\ \ .
\]
Another is
\[%
% [tikz figure removed]
%BeginExpansion
% [tikz figure removed]%
%EndExpansion
\ \ .
\]
(When drawing cycle digraphs, one commonly tries to place the nodes in such a
way as to make the arcs as short as possible. Thus, it is natural to keep the
cycles separate in the picture. But formally speaking, any picture is fine, as
long as the nodes and arcs don't overlap.)
\end{example}

\begin{example}
Let $\sigma:\left[  10\right]  \rightarrow\left[  10\right]  $ be the map that
sends the elements $1,2,3,4,5,6,7,8,9,10$ to $5,4,3,2,6,10,1,9,8,7$,
respectively. Then, $\sigma$ is a bijection, hence a permutation of $\left[
10\right]  $. The one-line notation of $\sigma$ is $\left(
5,4,3,2,6,10,1,9,8,7\right)  $. If we omit the commas and the parentheses,
then this becomes%
\[
5\ 4\ 3\ 2\ 6\ \left(  10\right)  \ 1\ 9\ 8\ 7.
\]
(We have put the $10$ in parentheses to make its place clearer.) The cycle
digraph of $\sigma$ is%
\[%
% [tikz figure removed]
%BeginExpansion
% [tikz figure removed]%
%EndExpansion
\ \ .
\]

\end{example}

Note that two-line notations and cycle digraphs can be defined (just as we did
in Definition \ref{def.perm.notations}) not just for permutations $\sigma\in
S_{n}$, but more generally for permutations $\sigma\in S_{X}$ of any finite
set $X$. But the one-line notation makes sense for $\sigma\in S_{n}$ only.


\subsection{Transpositions, cycles and involutions}

We shall now define some important families of permutations.

\subsubsection{Transpositions}

\begin{definition}
\label{def.perm.tij}Let $i$ and $j$ be two distinct elements of a set $X$.

Then, the \emph{transposition} $t_{i,j}$ is the permutation of $X$ that sends
$i$ to $j$, sends $j$ to $i$, and leaves all other elements of $X$ unchanged.
\end{definition}

Strictly speaking, the notation $t_{i,j}$ is somewhat ambiguous, since it
suppresses $X$. However, most of the times we will use it, the set $X$ will be
either clear from the context or irrelevant.

\begin{example}
The permutation $t_{2,4}$ of the set $\left[  7\right]  $ sends the elements
$1,2,3,4,5,6,7$ to $1,4,3,2,5,6,7$, respectively. Its one-line notation (with
commas and parentheses omitted) is therefore $1432567$.
\end{example}

Note that $t_{i,j}=t_{j,i}$ for any two distinct elements $i$ and $j$ of a set
$X$.

\begin{definition}
\label{def.perm.si}Let $n\in\mathbb{N}$ and $i\in\left[  n-1\right]  $. Then,
the \emph{simple transposition} $s_{i}$ is defined by%
\[
s_{i}:=t_{i,i+1}\in S_{n}.
\]

\end{definition}

Thus, a simple transposition is a transposition that swaps two consecutive
integers. Again, the notation $s_{i}$ suppresses $n$, but this won't usually
be a problem.

\begin{example}
The permutation $s_{2}$ of the set $\left[  7\right]  $ sends the elements
$1,2,3,4,5,6,7$ to $1,3,2,4,5,6,7$, respectively. Its one-line notation is
therefore $1324567$.
\end{example}

Here are some very basic properties of simple transpositions:\footnote{Recall
Definition \ref{def.perm.perm}. Thus, for example, $s_{i}^{2}$ means
$s_{i}s_{i}=s_{i}\circ s_{i}$, whereas $s_{i}s_{j}$ means $s_{i}\circ s_{j}$.}

\begin{proposition}
\label{prop.perm.si.rules}Let $n\in\mathbb{N}$. \medskip

\textbf{(a)} We have $s_{i}^{2}=\operatorname*{id}$ for all $i\in\left[
n-1\right]  $. In other words, we have $s_{i}=s_{i}^{-1}$ for all $i\in\left[
n-1\right]  $. \medskip

\textbf{(b)} We have $s_{i}s_{j}=s_{j}s_{i}$ for any $i,j\in\left[
n-1\right]  $ with $\left\vert i-j\right\vert >1$. \medskip

\textbf{(c)} We have $s_{i}s_{i+1}s_{i}=s_{i+1}s_{i}s_{i+1}$ for any
$i\in\left[  n-2\right]  $.
\end{proposition}

\begin{proof}
To prove that two permutations $\alpha$ and $\beta$ of $\left[  n\right]  $
are identical, it suffices to show that $\alpha\left(  k\right)  =\beta\left(
k\right)  $ for each $k\in\left[  n\right]  $. Using this strategy, we can
prove all three parts of Proposition \ref{prop.perm.si.rules}
straightforwardly (distinguishing cases corresponding to the relative
positions of $k$, $i$, $i+1$, $j$ and $j+1$). This is done in detail for
Proposition \ref{prop.perm.si.rules} \textbf{(c)} in \cite[solution to
Exercise 5.1 \textbf{(a)}]{detnotes}; the proofs of parts \textbf{(a)} and
\textbf{(b)} are easier and LTTR.
\end{proof}

\subsubsection{Cycles}

The following definition can be viewed as a generalization of Definition
\ref{def.perm.tij}:

\begin{definition}
\label{def.perm.cycs}Let $X$ be a set. Let $i_{1},i_{2},\ldots,i_{k}$ be $k$
distinct elements of $X$. Then,%
\[
\operatorname*{cyc}\nolimits_{i_{1},i_{2},\ldots,i_{k}}%
\]
means the permutation of $X$ that sends%
\begin{align*}
&  i_{1}\text{ to }i_{2},\\
&  i_{2}\text{ to }i_{3},\\
&  i_{3}\text{ to }i_{4},\\
&  \ldots,\\
&  i_{k-1}\text{ to }i_{k},\\
&  i_{k}\text{ to }i_{1}%
\end{align*}
and leaves all other elements of $X$ unchanged. In other words,
$\operatorname*{cyc}\nolimits_{i_{1},i_{2},\ldots,i_{k}}$ means the
permutation of $X$ that satisfies%
\begin{align*}
\operatorname*{cyc}\nolimits_{i_{1},i_{2},\ldots,i_{k}}\left(  p\right)   &  =%
\begin{cases}
i_{j+1}, & \text{if }p=i_{j}\text{ for some }j\in\left\{  1,2,\ldots
,k\right\}  ;\\
p, & \text{otherwise}%
\end{cases}
\\
&
\ \ \ \ \ \ \ \ \ \ \ \ \ \ \ \ \ \ \ \ \ \ \ \ \ \ \ \ \ \ \ \ \ \ \ \ \ \ \ \ \text{for
every }p\in X,
\end{align*}
where $i_{k+1}$ means $i_{1}$.

This permutation is called a $k$\emph{-cycle}.
\end{definition}

The name \textquotedblleft$k$-cycle\textquotedblright\ harkens back to the
cycle digraph of $\operatorname*{cyc}\nolimits_{i_{1},i_{2},\ldots,i_{k}}$,
which consists of a cycle of length $k$ (containing the nodes $i_{1}%
,i_{2},\ldots,i_{k}$ in this order) along with $\left\vert X\right\vert -k$
isolated nodes (more precisely, each of the elements of $X\setminus\left\{
i_{1},i_{2},\ldots,i_{k}\right\}  $ has an arrow from itself to itself in the
cycle digraph of $\sigma$). Here is an example:

\begin{example}
Let $X=\left[  8\right]  $. Then, the permutation $\operatorname*{cyc}%
\nolimits_{2,6,5}$ of $X$ sends
\begin{align*}
&  2\text{ to }6,\\
&  6\text{ to }5,\\
&  5\text{ to }2
\end{align*}
and leaves all other elements of $X$ unchanged. Thus, this permutation has OLN
$16342578$ and cycle digraph
\[%
% [tikz figure removed]
%BeginExpansion
% [tikz figure removed]%
%EndExpansion
\]

\end{example}

\begin{example}
Let $X$ be a set. If $i$ and $j$ are two distinct elements of $X$, then
$\operatorname*{cyc}\nolimits_{i,j}=t_{i,j}$. Thus, the $2$-cycles in $S_{X}$
are precisely the transpositions in $S_{X}$, so there are $\dbinom{\left\vert
X\right\vert }{2}$ many of them (since any $2$-element subset $\left\{
i,j\right\}  $ of $X$ gives rise to a transposition $t_{i,j}$, and this
assignment of transpositions to $2$-element subsets is bijective).
\end{example}

Note that the $k$-cycle $\operatorname*{cyc}\nolimits_{i_{1},i_{2}%
,\ldots,i_{k}}$ is often denoted by $\left(  i_{1},i_{2},\ldots,i_{k}\right)
$, but I will not use this notation here, since it clashes with the standard
notation for $k$-tuples.

\begin{exercise}
\label{exe.perm.cyc.how-many-kcyc}Let $n\in\mathbb{N}$ and let $k\in\left[
n\right]  $. Let $X$ be an $n$-element set. How many $k$-cycles exist in
$S_{X}$ ?
\end{exercise}

\begin{proof}
[Solution to Exercise \ref{exe.perm.cyc.how-many-kcyc} (sketched).]First, we
note that there is exactly one $1$-cycle in $S_{X}$ (for $n>0$), since a
$1$-cycle is just the identity map. This should be viewed as a degenerate
case; thus, we WLOG assume that $k>1$.

For any $k$ distinct elements $i_{1},i_{2},\ldots,i_{k}$ of $X$, we have%
\[
\operatorname*{cyc}\nolimits_{i_{1},i_{2},\ldots,i_{k}}=\operatorname*{cyc}%
\nolimits_{i_{2},i_{3},\ldots,i_{k},i_{1}}=\operatorname*{cyc}\nolimits_{i_{3}%
,i_{4},\ldots,i_{k},i_{1},i_{2}}=\cdots=\operatorname*{cyc}\nolimits_{i_{k}%
,i_{1},i_{2},\ldots,i_{k-1}}.
\]
That is, $\operatorname*{cyc}\nolimits_{i_{1},i_{2},\ldots,i_{k}}$ does not
change if we cyclically rotate the list $\left(  i_{1},i_{2},\ldots
,i_{k}\right)  $.

Any $k$-cycle $\operatorname*{cyc}\nolimits_{i_{1},i_{2},\ldots,i_{k}}$
uniquely determines the elements $i_{1},i_{2},\ldots,i_{k}$ up to cyclic
rotation (since $k>1$). Indeed, if $\sigma=\operatorname*{cyc}\nolimits_{i_{1}%
,i_{2},\ldots,i_{k}}$ is a $k$-cycle, then the elements $i_{1},i_{2}%
,\ldots,i_{k}$ are precisely the elements of $X$ that are not fixed by
$\sigma$ (it is here that we use our assumption $k>1$), and furthermore, if we
know which of these elements is $i_{1}$, then we can reconstruct the remaining
elements $i_{2},i_{3},\ldots,i_{k}$ recursively by%
\[
i_{2}=\sigma\left(  i_{1}\right)  ,\ \ \ \ \ \ \ \ \ \ i_{3}=\sigma\left(
i_{2}\right)  ,\ \ \ \ \ \ \ \ \ \ i_{4}=\sigma\left(  i_{3}\right)
,\ \ \ \ \ \ \ \ \ \ \ldots,\ \ \ \ \ \ \ \ \ \ i_{k}=\sigma\left(
i_{k-1}\right)
\]
(that is, $i_{2},i_{3},\ldots,i_{k}$ are obtained by iteratively applying
$\sigma$ to $i_{1}$). Therefore, if $\sigma\in S_{X}$ is a $k$-cycle, then
there are precisely $k$ lists $\left(  i_{1},i_{2},\ldots,i_{k}\right)  $ for
which $\sigma=\operatorname*{cyc}\nolimits_{i_{1},i_{2},\ldots,i_{k}}$ (coming
from the $k$ possibilities for which of the $k$ non-fixed points of $\sigma$
should be $i_{1}$).

Hence, the map%
\begin{align*}
f:\left\{  k\text{-tuples of distinct elements of }X\right\}   &
\rightarrow\left\{  k\text{-cycles in }S_{X}\right\}  ,\\
\left(  i_{1},i_{2},\ldots,i_{k}\right)   &  \mapsto\operatorname*{cyc}%
\nolimits_{i_{1},i_{2},\ldots,i_{k}}%
\end{align*}
is a $k$-to-$1$ map (i.e., each $k$-cycle in $S_{X}$ has precisely $k$
preimages under this map). Therefore, Lemma \ref{lem.count.multijection}
(applied to $m=k$ and \newline$A=\left\{  k\text{-tuples of distinct elements
of }X\right\}  $ and $B=\left\{  k\text{-cycles in }S_{X}\right\}  $) yields%
\[
\left(  \text{\# of }k\text{-tuples of distinct elements of }X\right)
=k\cdot\left(  \text{\# of }k\text{-cycles in }S_{X}\right)  .
\]
Therefore,
\begin{align*}
\left(  \text{\# of }k\text{-cycles in }S_{X}\right)   &  =\dfrac{1}{k}%
\cdot\underbrace{\left(  \text{\# of }k\text{-tuples of distinct elements of
}X\right)  }_{\substack{=n\left(  n-1\right)  \left(  n-2\right)
\cdots\left(  n-k+1\right)  \\\text{(by (\ref{eq.count.k-tups-dist-elts}),
since }X\text{ is an }n\text{-element set)}}}\\
&  =\dfrac{1}{k}\cdot n\left(  n-1\right)  \left(  n-2\right)  \cdots\left(
n-k+1\right) \\
&  =\dbinom{n}{k}\cdot\left(  k-1\right)  !\ \ \ \ \ \ \ \ \ \ \left(
\text{by a bit of simple algebra}\right)  .
\end{align*}
This is the answer to Exercise \ref{exe.perm.cyc.how-many-kcyc} in the case
$k>1$. Hence, Exercise \ref{exe.perm.cyc.how-many-kcyc} is solved.
\end{proof}

\subsubsection{Involutions}

\emph{Involutions} are maps that are inverse to themselves:

\begin{definition}
\label{def.perm.invol}Let $X$ be a set. An \emph{involution} of $X$ means a
map $f:X\rightarrow X$ that satisfies $f\circ f=\operatorname*{id}$. Clearly,
an involution is always a permutation, and equals its own inverse.
\end{definition}

For example, the identity map $\operatorname*{id}\nolimits_{X}$ is an
involution, and any transposition $t_{i,j}\in S_{X}$ is an involution, whereas
a $k$-cycle $\operatorname*{cyc}\nolimits_{i_{1},i_{2},\ldots,i_{k}}$ with
$k>2$ is never an involution. Further examples of involutions are products of
disjoint transpositions -- i.e., permutations of the form $t_{i_{1},j_{1}%
}t_{i_{2},j_{2}}\cdots t_{i_{k},j_{k}}$, where all $2k$ elements $i_{1}%
,j_{1},i_{2},j_{2},\ldots,i_{k},j_{k}$ are distinct. In fact, any involution
of a finite set can be expressed in this form.

An involution $\sigma\in S_{n}$ has a cycle digraph consisting of $1$-cycles
and $2$-cycles. For example, here is the cycle digraph of the involution
$t_{1,5}t_{2,3}\in S_{7}$:%
\[%
% [tikz figure removed]
%BeginExpansion
% [tikz figure removed]%
%EndExpansion
\]

